
\section{Results}

\subsection{Quality of Evidence Collected through \name{}}
To answer \textbf{RQ1}, we evaluate the quality of the \quantempplus{} dataset both qualitatively and quantitatively.

To \textbf{quantitatively} measure the efficacy of our dataset, we fine-tune \robertamnlilarge{} with the training split of our dataset and evaluated its performance on our test split of claim and evidence pairs collecting using queries from \name{} based decomposition. From Table.\ref{tab:data_quality_quantitative} we see that the model trained on our dataset, \robertaqtempplus{}, performs much better than all the baselines on all metrics. It achieves relative gains of about 41\% over \naive{} and about 16\% over the \robertaclaimnoev{} classifier. We also see that the performance of \robertaqtempplus{} is closer to \robertagold{} (uses justification documents directly) indicating that we are very close to the verification upper bound. This demonstrates that queries from \emph{\name{} which emulates the search process in justification snippets from fact-checkers, are of better quality, resulting in better quality relevant evidences for each claim}. It is important to note \robertaqtempplus{} uses evidence without temporal leakage or gold justification leakage.

We also perform the qualitative evaluation of the queries from \name{} and resulting evidence collection in \quantempplus{} benchmark. Following the setup of ClaimDecomp \cite{claimdecomp}, we sample 50 representative claims across all categories of numerical claims to perform this analysis with help of 3 annotators who are researchers working on fact-checking. The inter-annotator agreement in presented in Table \ref{table:agreement}. We observe that the average precision of these queries is 90\%, meaning most of our queries are \textbf{relevant} to the claim. Additionally, we observe an average completeness score of \textbf{4.49} indicating queries from \name{} cover most of the aspects of the claim. The average redundancy score is \textbf{2.4} indicating mild redundancy due to our filtering process. Likewise for collected evidence snippets we observe an average completeness of \textbf{3.94} and redundancy of \textbf{2.9}.

\begin{table}[hbt!]
\centering
\begin{tabular}{lcc}
\toprule
\textbf{Metric} & \textbf{Queries} & \textbf{Evidence Snippets} \\
\midrule
\textbf{Completeness} & \cellcolor{green!10}0.61 & \cellcolor{green!10}0.66 \\
\textbf{Redundancy} & \cellcolor{yellow!10}0.4 & \cellcolor{yellow!10}0.45 \\
\textbf{Relevance} & \cellcolor{green!10}0.58 & \cellcolor{green!10}0.67 \\
\bottomrule
\end{tabular}
\caption{Inter-annotator agreement (Fleiss’) for qualitative analysis }%\quantempplus{}}
\label{table:agreement}
\vspace{-4em}
\end{table}
%\vspace{-2em}
\subsection{Retrieval Performance on \quantempplus{} }

\label{5.2}



To answer \textbf{RQ2}, we compare the retrieval performance of various claim decomposition approaches as shown in Table \ref{tab:ret_perf_temp}.  The retrieval is performed in an open-domain setup from entire evidence collection which also contains distractors. Over the top-k documents retrieved from the corpus using Contriever, we also add a \textbf{temporal filter} to retain only top-k documents that are published before the claim publication date. Though we already employ the temporal filter per-claim level when collecting evidence from the web, temporal leakage can still occur because evidences collected for other related claims in the benchmark maybe published at a later date. Hence, in an open-domain retrieval setup such evidences which are semantically related could be retrieved leading to temporal leakage. Prior benchmarks like AveriTec \cite{averitec} \textbf{do not account for this form of leakage in open-domain setup}. We observe that \name{} with temporal filter achieves the highest Recall@k for k=10,100 and also highest MRR which indicates higher ranking quality. This is primarily because \name{} generates queries by extracting the manual fact-checker's search process from justification documents.
\begin{table}[ht]
    \centering
    \scalebox{1.0}{
    \begin{tabular} { 
      lcccc}
        \toprule
        \textbf{Query mode} & \textbf{NDCG@10} & \textbf{Recall@10} & \textbf{Recall@100}  & \textbf{MRR} \\ \midrule
        \claimonly{} & 0.33 & 0.32 & 0.51 &  0.54 \\ \midrule
        \textbf{Decomposition}& & &  &\\ \midrule
        \oracleq{} & 0.57 & 0.54 & 0.82 & 0.68 \\ 
             \oracleq{}   (-temporal filter) & 0.54 & 0.51 & 0.82 & 0.651 \\
        \claimdecomp{} & 0.34 & 0.32 & 0.57  & 0.53 \\ 

                \claimdecomp{} (-temporal filter) & 0.31 & 0.30 & 0.56  & 0.49 \\ 
        \pgmfc{} & 0.32 & 0.31 & 0.52  &0.52\\ 

                \pgmfc{} (-temporal filter) & 0.30 & 0.29 & 0.52  &0.481\\ 
        \qgen{} & 0.33 & 0.31 & 0.56 & 0.57\\
          \qgen{} (-temporal filter) & 0.29 & 0.27 & 0.54 & 0.48\\\bottomrule
    \end{tabular}
    }
    \caption{Retrieval performance of different query planning methods using \combmaxn{} on \quantempplus{} with temporal filtering.}
    \label{tab:ret_perf_temp}
    \vspace{-2em}
\end{table}
While retrieval happens in a realistic setup one cannot assume access to justification documents as background information at test time to generate queries. Hence \name{} presents an upper bound on quality of queries and their impact on retrieval.  We observe that other decomposition approaches fall short in retrieval quality compared to \name{} as they try to decompose the claim without background information. Even \qgen{} which tries to distill \name{} queries from a LLM does not obtain optimal performance due to generation of queries based on claim alone without additional context. Hence, we hypothesize that iterative approaches that  claim decomposition and retrieval are necessary for optimal claim decomposition. We believe that by using additional information retrieved adaptively using intermediate queries/sub-claims to generate further queries  would help achieve better queries like \name{} and gains in retrieval performance. One \textbf{future direction}, is \emph{to develop new claim decomposition approaches  using queries from \name{} as ground truth}.

\vspace{-1em}

\subsection{Downstream Fact-verification Performance}
\label{5.3}
%\vspace{-1em}
While the above results indicate the performance of the query planning methods at retrieval, observing their corresponding downstream impact is essential as it represents the final outcome of the pipeline. Hence, to answer \textbf{RQ3} we evaluate the final performance of the veracity prediction component for each of the query planning methods, the results of which are shown in Table.\ref{table:downstream_nli}. 



\emph{Superiority of decomposition-based methods.} Firstly, we observe that the performance of \robertaclaimnoev{} performs the worst as it relies only on surface level patterns in claim for veracity prediction without evidence. Secondly, we observe that the downstream performance provided by decomposition-based query planning methods, specifically that of \qgen{} is higher than that of \claimonly{}. 

\begin{table}[h!]
\centering
\begin{tabular}{lcccccc}
\toprule
\textbf{Query mode} & \multicolumn{6}{c}{\textbf{NLI performance}} \\
\cmidrule(lr){2-7}
& \textbf{Accuracy} & \textbf{F1-T}&\textbf{F1-F}&\textbf{F1-C}&\textbf{W-F1} & \textbf{M-F1} \\
\midrule
\robertaclaimnoev{} & 63.2 &24.71&79.93&47.37 & 61.53 & 50.67 \\
\oracle{} & 65.25&47.76&79.92&49.67 & 66.56 & 59.11 \\
\claimonly{} & 66.53&53.38&81.46&31.02 & 64.03 & \textbf{55.28} \\ \midrule
\textbf{Decomposition} & & &  \\ \midrule
\oracleq{} & 66.45&45.73&80.78&50.25 & 66.8 & \textbf{58.92}\textsuperscript{\textdagger} \\
\claimdecomp{} & 64.57&53.78&79.97&35.81 & 64.41 & 56.51 \\
\pgmfc{} & 65.67&51.38&81.19&37.15 & 65.73 & 56.57 \\
\qgen{} & 66.25&53.62&81.39&36.95 & 65.41 & \textbf{57.28}\textsuperscript{\textdagger} \\

\bottomrule
\end{tabular}
\caption{Fact Verification performance using different claim decomposition methods on \quantempplus{}.\textsuperscript{\textdagger} indicates statistical significance over \claimonly{} at 0.05 level. F1-T, F1-F, and F1-C - True, False, and Conflicting F1 scores.}
\label{table:downstream_nli}
    \vspace{-2em}
\end{table}

While \pgmfc{} and \claimdecomp{} which employ a larger model \gptt{} for decomposition are minimally better than \claimonly{}, \qgen{} which employs a much smaller model (Flan-T5-Large) provides a statistically significant gain of 3.6\% in the Macro-F1 score over \claimonly{} and is marginally better than \pgmfc{} and \claimdecomp{}. A comparison of the per-class F1 scores shows that decomposition-based query planning methods excel at claims that are of a conflicting nature. We observe relative gains of up to 20\% by \pgmfc{} compared to \claimonly{}. Claims whose veracity is of conflicting nature specifically require diverse perspectives to be retrieved as evidence since some parts of them may be true and other parts false. We observe that \claimonly{} fails to provide the correct downstream result as the evidence it retrieves depends only on the claim and is too homogeneous failing to capture multiple explicit and implicit aspects.


% \emph{Performance of \qgen{}.} Another interesting observation we see is that, despite \pgmfc{} and \claimdecomp{} using a larger LLM \gptt{} for decomposition, the downstream impact of \qgen{}'s claim decomposition method, which was developed by training a smaller model of Flan-T5 with our oracle queries from \quantempplus{}, is superior. The training of \qgen{} can also be seen as a form of knowledge distillation \cite{alpaca,vicuna,peng2023instruction,gu2024minillm}, as our oracle queries were generated by prompting the \gptt{} model with the claim and justification document as input.  

%Additionally, \cite{divide-conquer} has shown that problem decomposition tasks are easier to distill into small LLMs in QA tasks. We observe the same with decomposing claims into sub-queries, which aids in retrieving and verifying information.

\textbf{Discrepancy between retrieval and downstream performance:} An interesting observation that we see from comparing the results from retrieval in Table.\ref{tab:ret_perf_temp} and NLI from Table.\ref{table:downstream_nli} is that the gains in retrieval performance do not proportionally translate to downstream gains in NLI (fact-verification) performance. The retrieval performance of queries from \name{} provides significant relative gains %of %68.7\% at Recall@10 and 74\% at NDCG@10
over using \claimonly{} for retrieval. However, we see that this only translated to a relative gain of 6.6\% downstream. We observe that this may be \emph{primarily due to fundamental  numerical contextualization and understanding limitations in claim verification models and LLMs}. These models lack abilities to parse, contextualize, understand notion of ranges, thresholds and other numerical information in presence of other textual information.  Additionally, while retrieval systems optimize for end user consumption \cite{ret-enh-ml}, the retrieved evidences might not necessarily be useful for downstream tasks. Inspired by prior works in other tasks \cite{zhang-etal-2023-relevance} one could train the claim decomposition model by optimizing for downstream fact-verification performance in an end-end manner to reduce the observed gap in the performance of retrieval and downstream NLI. 



